%%%%%%%%%%%%%%%%%%%%%%%%%%%%%%%%%%%%%%%%%%%%%%%%%%%%%%%%%%%%%%%%%%%%%%
% Lớp 11 | Hàm số lượng giác và phương trình lượng giác | [Giảm] Phương trình lượng giác thường gặp
% Dạng 1D1?6-6: Phương trình lượng giác có chứa ẩn ở mẫu số
%%%%%%%%%%%%%%%%%%%%%%%%%%%%%%%%%%%%%%%%%%%%%%%%%%%%%%%%%%%%%%%%%%%%%%


% nguồn: 2026-11-15_11_1D1_Luong_giac (nhập 14/11/2026)
\begin{bt}%[1D1V6-6]%[ma:31f5be9b]
% nguon: Archive/SP_KhoLatex/ID6_old/1-CK1-OT-De-18.tex #2827

Giải phương trình $\dfrac{1}{\cos x}+\dfrac{\sqrt{3}}{\sin x}=2 \sin \left(x+\dfrac{\pi}{3}\right)$.
	\loigiai{
		Điều kiện $\heva{&\sin x \neq 0 \\ &\cos x \neq 0} \Leftrightarrow \sin 2 x \neq 0 \Leftrightarrow 2 x \neq k \pi \Leftrightarrow x \neq \dfrac{k \pi}{2} \quad(k \in \mathbb{Z})$.\\
		Theo đề bài ta có
		\begin{eqnarray*}
			&& \dfrac{1}{\cos x}+\dfrac{\sqrt{3}}{\sin x}=2 \sin \left(x+\dfrac{\pi}{3}\right) \\
			& \Leftrightarrow& \dfrac{\sin x+\sqrt{3} \cos x}{\sin x \cos x}=2 \sin \left(x+\dfrac{\pi}{3}\right) \\
			&\Leftrightarrow& \dfrac{\dfrac{1}{2} \sin x+\dfrac{\sqrt{3}}{2} \cos x}{\sin x \cos x}=\sin \left(x+\dfrac{\pi}{3}\right) \\
			& \Leftrightarrow& \dfrac{\sin \left(x+\dfrac{\pi}{3}\right)}{\sin x \cos x}=\sin \left(x+\dfrac{\pi}{3}\right) \\
			 &\Leftrightarrow& \sin \left(x+\dfrac{\pi}{3}\right)\left(\dfrac{1}{\sin x \cos x}-1\right)=0\\
	&	\Leftrightarrow& \sin \left(x+\dfrac{\pi}{3}\right)\left(\dfrac{2}{\sin 2 x}-1\right)=0 \\
		& \Leftrightarrow&\hoac{&
			\sin \left(x+\dfrac{\pi}{3}\right)=0 \\
			&\sin 2 x=2	} 
		\end{eqnarray*}
	Với $\sin \left(x+\dfrac{\pi}{3}\right)=0 \Leftrightarrow x+\dfrac{\pi}{3}=k \pi \Leftrightarrow x=-\dfrac{\pi}{3}+k \pi \quad(k \in \mathbb{Z}).$\\		
		Với $\sin 2 x=2$ phương trình vô nghiệm.\\
		Kết hợp với điều kiện, vậy nghiệm của phương trình là $x=-\dfrac{\pi}{3}+k \pi \quad(k \in \mathbb{Z})$.
	}
\end{bt}

% nguồn: 2026-11-15_11_1D1_Luong_giac (nhập 14/11/2026)
\begin{bt}%[1D1C6-6]%[ma:ad340e6f]
% nguon: Archive/SP_KhoLatex/ID6_old/1-GK1-KNTT-De16-NH23-24.tex #3038

Giải phương trình: $\cot x+\sin x\left(1+\tan x \tan \dfrac{x}{2}\right)=4$.
	\loigiai{
	Điều kiện: $
	\heva{&\sin x \neq 0\\&\cos x \neq 0\\&\cos \dfrac{x}{2} \neq 0}
	\Leftrightarrow
	\heva{&\sin x \neq 0\\&\cos x \neq 0}
    \Leftrightarrow \sin 2 x \neq 0$.
    	\begin{eqnarray*}
    	& & \cot x+\sin x\left(1+\tan x \tan \dfrac{x}{2}\right)=4\\ 
    	& \Leftrightarrow & \dfrac{\cos x}{\sin x}+\sin x\left(1+\dfrac{\sin x}{\cos x} \cdot \dfrac{\sin \dfrac{x}{2}}{\cos \dfrac{x}{2}}\right)=4\\
    	& \Leftrightarrow &  \dfrac{\cos x}{\sin x}+\sin x\left(\dfrac{\cos x \cos \dfrac{x}{2}+\sin x \cdot \sin \dfrac{x}{2}}{\cos x \cdot \cos \dfrac{x}{2}}\right)=4\\
    	& \Leftrightarrow & \dfrac{\cos x}{\sin x}+\sin x \cdot \dfrac{\cos \dfrac{x}{2}}{\cos x \cdot \cos \dfrac{x}{2}}=4\\ 
    	& \Leftrightarrow & \dfrac{\cos x}{\sin x}+\dfrac{\sin x}{\cos x}=4\\
    	& \Leftrightarrow & \dfrac{\cos ^2 x+\sin ^2 x}{\sin x \cdot \cos x}=4\\
        & \Leftrightarrow & 4 \sin x \cdot \cos x=1\\
        & \Leftrightarrow & 2 \sin 2 x=1 \Leftrightarrow \sin 2 x=\dfrac{1}{2}\\
        & \Leftrightarrow & \hoac{&x=\dfrac{\pi}{12}+k \pi\\&x=\dfrac{5 \pi}{12}+k \pi}, k \in \mathbb{Z}.
    \end{eqnarray*}	
So với điều kiện, ta nhận 
$\hoac{&x=\dfrac{\pi}{12}+k \pi \\&x=\dfrac{5 \pi}{12}+k \pi}, k \in \mathbb{Z}$.
	}
\end{bt}

% nguồn: 2026-11-15_11_1D1_Luong_giac (nhập 14/11/2026)
\begin{bt}%[1D1H6-6]%[ma:c1132a4a]
% nguon: Archive/SP_KhoLatex/ID6_old/11-OT-PTLuongGiacCoBan.tex #3096
%[ĐGNL-ĐGTD-2022, Tran Tony]
	Giải phương trình $\tan 3x\cdot\cot 2x=1$.
	\loigiai{
		Điều kiện: $\heva{&\cos 3x\neq 0 \\ &\sin 2x\neq 0}\Leftrightarrow\heva{&x\neq \dfrac{\pi }{6}+k\dfrac{\pi }{3} \\ &x\neq k\dfrac{\pi }{2}}, k\in \mathbb{Z}$.\\
		Phương trình $\Leftrightarrow \tan 3x=\dfrac{1}{\cot 2x}=\tan 2x$ $\Leftrightarrow 3x=2x+k\pi $
		$\Leftrightarrow x=k\pi$, $\ k\in \mathbb{Z}$ (nghiệm này bị loại).\\
		Vậy phương trình đã cho vô nghiệm.
	}
\end{bt}

% nguồn: 2026-11-15_11_1D1_Luong_giac (nhập 14/11/2026)
\begin{bt}%[1D1V6-6]%[ma:8ff5b927]
% nguon: Archive/SP_KhoLatex/ID6_old/12-OT-PTLuongGiacThuongGap.tex #3698

	Giải phương trình $\dfrac{\sqrt{3}\sin x-\cos x-1}{\sin 2x}=0$.
	\loigiai{
		Điều kiện $\sin 2x\ne 0\Leftrightarrow x\ne \dfrac{k\pi}{2}$ với $k\in\mathbb{Z}$.\\
		Khi đó phương trình đã cho trở thành
		\allowdisplaybreaks
		\begin{eqnarray*}
			\sqrt{3}\sin x-\cos x-1=0&\Leftrightarrow&\dfrac{\sqrt{3}}{2}\sin x-\dfrac{1}{2}\cos x=\dfrac{1}{2}\\
			&\Leftrightarrow&\sin\left(x-\dfrac{\pi}{6}\right)=\dfrac{1}{2}\\
			&\Leftrightarrow&\hoac{&x-\dfrac{\pi}{6}=\dfrac{\pi}{6}+k2\pi\\&x-\dfrac{\pi}{6}=\dfrac{5\pi}{6}+k2\pi}\\
			&\Leftrightarrow&\hoac{&x=\dfrac{\pi}{3}+k2\pi\\&x=\pi+k2\pi}, \left(k\in\mathbb{Z}\right).
		\end{eqnarray*}
		So sánh với điều kiện xác định của phương trình ta có nghiệm của phương trình là $x=\dfrac{\pi}{3}+k2\pi$ với $k\in\mathbb{Z}$.
	}
\end{bt}

% nguồn: 2026-11-15_11_1D1_Luong_giac (nhập 14/11/2026)
\begin{bt}%[1D1H6-6]%[ma:7ee41195]
% nguon: Archive/SP_KhoLatex/ID6_old/12-OT-PTLuongGiacThuongGap.tex #3705

	Giải phương trình $\dfrac{\cos x+ \sin 2x}{\cos 3x}+1=0$.
	\loigiai{
	\begin{itemize} 
 \item Điều kiện $ \cos 3x \ne 0 \Leftrightarrow x\ne \dfrac{\pi}{6}+\dfrac{k \pi }{3},\,\,(k\in \mathbb{Z})$.
 \item Phương trình đã cho tương đương với
 \begin{eqnarray*} 
 	& & \cos x-\cos 3x +\sin 2x=0\\ 
 	&\Leftrightarrow & 2\sin 2x \sin x +\sin 2x=0\\ 
 	&\Leftrightarrow & \hoac{&\sin 2x =0\\&\sin x=-\dfrac{1}{2}} \Leftrightarrow \hoac{&x=\dfrac{k\pi}{2}\\&x=-\dfrac{\pi}{6}+k2\pi \\& x=\dfrac{7\pi}{6}+k2\pi}, \left(k\in\mathbb{Z}\right). 
 \end{eqnarray*} 
 Vì $\cos 3x=4\cos^3 x-3\cos x$ nên đối chiếu điều kiện ta phải loại đi những giá trị mà $\cos x=0$.\\
 Cuối cùng ta được $\hoac{&x=k\pi \\&x=-\dfrac{\pi}{6}+k2\pi \\& x=\dfrac{7\pi}{6}+k2\pi},\,\,(k\in \mathbb{Z})$.
 \end{itemize}
	} 
\end{bt}

% nguồn: chuyen_dang (nhập 22/11/2026)
\begin{bt}%[1D1V6-6]%[ma:631cdf8b]
% nguon: Archive/SP_KhoLatex/ID6_old/1K1-4-5.tex #3329
%[DCHT Toán 11 - KNTT -Thọ Bùi] 
	Giải phương trình $\sin^6x+\cos^6 x=\dfrac{7}{16}$.
	\loigiai{Ta có:
		\allowdisplaybreaks
		\begin{eqnarray*}
			\sin^6x+\cos^6 x=\dfrac{7}{16}
			&\Leftrightarrow&\left(\sin^2 x+\cos^2x\right)^3-3\sin^2x\cos^2x\left(\sin^2 x+\cos^2 x\right)=\dfrac{7}{16}\\
			&\Leftrightarrow&1-\dfrac{3}{4}\sin^2 2x=\dfrac{7}{16}\\
			&\Leftrightarrow&1-\dfrac{3}{4}\left(\dfrac{1-\cos 4x}{2}\right)=\dfrac{7}{16}\\
			&\Leftrightarrow&\cos 4x =-\dfrac{1}{2}\\
			&\Leftrightarrow&\cos 4x=\cos\dfrac{2\pi}{3}\\
			&\Leftrightarrow&x=\pm \dfrac{\pi}{6}+\dfrac{k\pi}{2}, \left(k\in \mathbb{Z}\right).
		\end{eqnarray*}
	}
\end{bt}

% nguồn: chuyen_dang (nhập 22/11/2026)
\begin{bt}%[1D1H6-6]%[ma:13c72c3a]
% nguon: Archive/SP_KhoLatex/ID6_old/GK1.6.tex #3330

	Giải phương trình $2\sin x-1=0$.
	\loigiai{
	Ta có:
	\begin{eqnarray*}
		&& 2\sin x-1=0\Leftrightarrow 2\sin x=1 \\
		& \Leftrightarrow& \sin x=\dfrac{1}{2}=\sin \dfrac{\pi}{6}\\
		&\Leftrightarrow& \hoac{& x=\dfrac{\pi}{6}+k2\pi \\
		& x=\pi-\dfrac{\pi}{6}+k2\pi} \\
		&\Leftrightarrow& \hoac{& x=\dfrac{\pi}{6}+k2\pi \\ & x=\dfrac{5\pi}{6}+k2\pi}(k\in \mathbb{Z}).
	\end{eqnarray*}
	}
\end{bt}

% nguồn: chuyen_dang (nhập 22/11/2026)
\begin{bt}%[1D1H6-6]%[ma:c8ebc69f]
% nguon: Archive/2024/TOAN_11_XE/TOAN_11_XE/data4/11C1/KNTT-1.4.tex #1088
%[DCHT Toán 11 - KNTT -Thọ Bùi] 
	Giải phương trình: $\sin x+\sin 2 x=0$.
	\loigiai{
		Ta có: \allowdisplaybreaks
		\begin{eqnarray*}
			\sin x+\sin 2 x=0&\Leftrightarrow& \sin x+2\sin x\cdot \cos x=0 \\
			&\Leftrightarrow& \sin x \cdot \left(1+2\cos x\right)=0 \\
			&\Leftrightarrow& \hoac{&\sin x=0\\&1+2\cos x=0} \\
			&\Leftrightarrow& \hoac{&\sin x=0\\&\cos x=-\dfrac{1}{2}}\Leftrightarrow \hoac{&x=k\pi\\&x=\pm \dfrac{2\pi}{3}+k2\pi} (k\in\mathbb{Z}).
		\end{eqnarray*}	
		Vậy phương trình có các nghiệm là $k\pi, k\in\mathbb{Z}$ và $x=\pm \dfrac{2\pi}{3}+k2\pi, k\in\mathbb{Z}$.
	}
\end{bt}

% nguồn: chuyen_dang (nhập 22/11/2026)
\begin{bt}%[1D1H6-6]%[ma:3c6efab6]
% nguon: Archive/SP_KhoLatex/ID6_old/11-OT-PTLuongGiacCoBan.tex #3093
%[ĐGNL-ĐGTD-2022, Tran Tony]
	Giải phương trình $\cos 3x = \dfrac{1}{3}$.
	\loigiai{Ta có  $\cos 3x = \dfrac{1}{3}\Leftrightarrow \hoac{&3x=\arccos\dfrac{1}{3}+k2\pi\\&3x=-\arccos\dfrac{1}{3}+k2\pi}\Leftrightarrow \hoac{&x=\dfrac{1}{3}\arccos\dfrac{1}{3}+\dfrac{k2\pi}{3}\\&x=\dfrac{1}{3}\arccos\dfrac{1}{3}+\dfrac{k2\pi}{3}}, (k\in \mathbb{Z}).$
	}	
\end{bt}

% nguồn: chuyen_dang (nhập 22/11/2026)
\begin{bt}%[1D1H6-6]%[ma:337ac33b]
% nguon: Archive/2024/TOAN_11_XE/TOAN_11_XE/data4/11C1/KNTT-1.4.tex #1075
%[Tex hóa SGK 11, KNTT]
	Giải phương trình $\sin 2x  = \sin (60^{\circ} - 3x)$.
	\loigiai{
Ta có 
\begin{align*}
\sin 2x = \sin (60^{\circ}-3x) & \Leftrightarrow \hoac{& 2x = 60^{\circ}-3x + k360^{\circ} \\ & 2x = 180^{\circ} - (60^{\circ}-3x) + k360^{\circ}}
		\Leftrightarrow \hoac{& 5x = 60^{\circ} + k360^{\circ} \\ & -x = 120^{\circ} + k360^{\circ}} \\
		&\Leftrightarrow \hoac{& x = 12^{\circ} + k72^{\circ} \\ & x = -120^{\circ} - k360^{\circ}} \quad (k \in \mathbb{Z}).
\end{align*}
	}
\end{bt}

% nguồn: chuyen_dang (nhập 22/11/2026)
\begin{bt}%[1D1V6-6]%[ma:ba22bf93]
% nguon: Archive/SP_KhoLatex/ID6_old/1K1-4-5.tex #3327
%[DCHT Toán 11 - KNTT -Thọ Bùi] 
	Giải phương trình: $\sin^4\dfrac{x}{2}+\cos^4 \dfrac{x}{2}=\dfrac{1}{2}$.
	\loigiai{ Ta có:
		\allowdisplaybreaks
		\begin{eqnarray*}
			\sin^4\dfrac{x}{2}+\cos^4 \dfrac{x}{2}=\dfrac{1}{2}
			&\Leftrightarrow&\left(\sin^2\dfrac{x}{2}+\cos^2\dfrac{x}{2}\right)^2-2\sin^2\dfrac{x}{2}\cos^2\dfrac{x}{2}=\dfrac{1}{2}\\
			&\Leftrightarrow&1-\dfrac{1}{2}\sin^2x=\dfrac{1}{2}\\
			&\Leftrightarrow&\sin^2 x=1\\
			&\Leftrightarrow&\cos x=0\\
			&\Leftrightarrow&x=\dfrac{\pi}{2}+k\pi, \left(k\in \mathbb{Z}\right).
		\end{eqnarray*}
		
	}
\end{bt}

% nguồn: chuyen_dang (nhập 22/11/2026)
\begin{ex}%[1D1H6-6]%[ma:592e4084]
	 Giải phương trình $9^{\sin 2x}=1$.
	 \loigiai{
		Ta có $9^{\sin 2x}=1\Leftrightarrow \sin  2x=0\Leftrightarrow 2x=k2\pi \Leftrightarrow x=k\pi, \left(k\in \mathbb{Z}\right)$.
	} 
\end{ex}

% nguồn: chuyen_dang (nhập 22/11/2026)
\begin{bt}%[1D1H6-6]%[ma:2a06d0bf]
% nguon: Archive/SP_KhoLatex/ID6_old/11-OT-PTLuongGiacCoBan.tex #3095
%[ĐGNL-ĐGTD-2022, Tran Tony]
	Giải phương trình $\sin x+\cos \left( x-\dfrac{\pi}{2}\right)=2$.
	\loigiai{
		$\sin x+\cos \left( x-\dfrac{\pi}{2}\right)=2\Leftrightarrow 2\sin x=2\Leftrightarrow \sin x=1\Leftrightarrow x=\dfrac{\pi}{2}+k2\pi,k\in\mathbb{Z}$.
	}
\end{bt}

% nguồn: chuyen_dang (nhập 22/11/2026)
\begin{bt}%[1D1H6-6]%[ma:3cb5ecfc]
% nguon: Archive/2024/TOAN_11_XE/TOAN_11_XE/data4/11C1/KNTT-1.4.tex #1085
%[DCHT Toán 11 - KNTT -Thọ Bùi] 
	Giải phương trình:  $\sin 4x - \cos\left(x + \dfrac{\pi}{6}\right)=0$.
	\loigiai{
		Ta có:
		\allowdisplaybreaks
		\begin{eqnarray*}
			&&\sin 4x - \cos\left(x + \dfrac{\pi}{6}\right)=0\\
			&\Leftrightarrow& \sin 4x = \cos\left(x + \dfrac{\pi}{6}\right)\\
			&\Leftrightarrow& \sin 4x = \sin \left(\dfrac{\pi}{3}-x\right)\\
			&\Leftrightarrow& \hoac{&4x=\dfrac{\pi}{3}-x+k2\pi\\&4x=\pi - \dfrac{\pi}{3}+x+k2\pi\\}\\
			&\Leftrightarrow& \hoac{&5x=\dfrac{\pi}{3}+k2\pi\\&3x= \dfrac{2\pi}{3}+k2\pi\\}\\
			&\Leftrightarrow& \hoac{&x=\dfrac{\pi}{15}+\dfrac{k2\pi}{5}\\&x= \dfrac{2\pi}{9}+\dfrac{k2\pi}{3}\\} (k\in \mathbb{Z}).
		\end{eqnarray*}
	}
\end{bt}
